\section{Identité structurelle et grandeur de masse
}
\label{vi:sec:introduccion}

Déterminer la masse d’une particule exige de préciser quel objet persiste et
quelle opération transforme sa structure en grandeur. Ce travail
construit ce lien : à partir des extensions compatibles et de leurs registres, on
obtient des fibres d’états, et sur celles-ci un opérateur positif de masse.
Sa spécialisation électronique réunit deux lecteurs coïncidant en une
même évaluation scalaire. La thèse est que l’identité, l’équation de
masse et la valeur évaluée appartiennent à une chaîne mathématique commune ; la
comparaison physique examine ensuite la réalisation de cette chaîne dans une
coordinatisation de mesure. Ainsi, les arguments et les coefficients de l’équation
conservent une provenance structurelle antérieure à la comparaison.


La dynamique de trajectoire incorpore l’histoire dans cette grandeur : elle détermine
le support cellulaire exact de la trajectoire centrale, reconstruit deux mémoires
qu’une statistique de comptage ne distingue pas, diagonalise le lecteur de retour
au moyen de caractères finis et caractérise la conservation de la masse par le
commutateur entre l’opérateur de masse et un transport TPK. Ces quatre
opérations agissent sur le même état et ne remplacent ni l’atlas, ni l’équation
électronique, ni leurs définitions et preuves antérieures.


On introduit un noyau formel mathématique dénommé holographie modulaire
triadique, ci-après HMT, composé de trois objets. \emph{All Possible
Paths}, ci-après APP, est la structure arithmétique de chemins sur
\((\mathbb Z/9\mathbb Z)^2\). L’adjacence est donnée par le graphe de
Cayley additif à générateurs cardinaux \(\pm(1,0)\) et \(\pm(0,1)\) ;
la réalisation cellulaire périodique de ce support est toroïdale. Ses
opérations d’accumulation additive et de composition
multiplicative se distinguent des évaluations qui déterminent leurs
coordonnées. TRIT est l’unité de signature orientée \(\{-1,0,+1\}\) :
elle détermine le régime local et l’orientation de l’opération. Le
\emph{Topological Phase Kernel}, ci-après TPK, constitue la dynamique
du noyau ; il compose sélection, transport, mise à jour de mémoire et
retour sur l’état enrichi. La mécanique dimensionnelle développe
les réalisations de cette structure en droites de grandeur et en espaces
d’états physiques. L’exposé construit ces objets à partir de leurs règles,
avant d’utiliser les coordonnées internes obtenues dans les opérateurs de masse.


Le développement obtient des résultats définis à chacun de ces niveaux. La
composition de \(108\) mises à jour du TPK produit un état terminal
qui conserve le registre incidentiel, ses deux canaux transversaux et la signature
dodécaphasique ; la transformée orbitale reconstruit de manière réversible le registre
\(K\). Sur cette généalogie sont construits l’atlas de
\(81/56/13/324\), l’opérateur positif de masse de chaque fibre, l’exception
scalaire de l’électron central et les règles de composition antérieures à
l’évaluation. La comparaison ultérieure maintient séparés valeurs centrales,
intervalles, limites, masses de pôle, masses courantes et largeurs. L’apport
consiste donc en une chaîne effective qui relie l’objet persistant,
sa mémoire de route, l’opérateur qui agit sur sa fibre et le type précis de
la grandeur évaluée.


Dans ce même ordre causal, le recensement complet
\(104\,976\to468\to243\subset729\) précède toute sélection régionale. La
bijection \(\beta:\mathbb F_3^6\to E_9^3\) identifie l’espace ambiant de
\(729\) positions, tandis que la différence de seuil produit la couronne
\(271=243+27+1\) et la complétion exacte
\(1000=729+243+27+1\). La mesure uniforme de la complétion est compatible
avec les projections dimensionnelles. Ces opérations construisent l’espace ambiant
et son bord ; la sélection ultérieure conserve son domaine propre.


\subsection{La structure qui précède le chiffre
}

La masse n’est pas la seule information d’une particule. Interviennent également
le support de persistance, ses orientations, la conjugaison, les
représentations et les règles de composition. Dans l’approche HMT,
ces données ne sont pas ajoutées à la fin à un chiffre isolé : elles appartiennent à la
construction sur laquelle agit la fonctionnelle de masse. Deux histoires peuvent
partager une lecture scalaire et conserver des informations différentes dans leurs
feuillets ou dans leur mémoire. C’est pourquoi l’égalité des valeurs doit être distinguée
de l’égalité des structures.


La distinction est particulièrement claire dans la composition. Un tableau peut
additionner des coordonnées déjà évaluées ; une dynamique doit d’abord déterminer si deux
histoires peuvent composer leur phase, leur bord, leurs orientations et leurs
prédécesseurs. Les registres se composent à ce niveau. La signature est extraite
ensuite, et le caractère évalue cette signature. La différence entre ces
opérations fournit le lieu mathématique permettant d’étudier les états composés,
les liaisons et les résonances.


La question de savoir pourquoi une grandeur vaut ce qu’elle vaut se trouve ainsi liée à
une autre plus précise : quels éléments de la construction déterminent sa valeur et
quelles informations restent disponibles après son évaluation. L’article
développe cette question à partir des règles internes du corpus. La
comparaison physique ultérieure examine ses réalisations sans intervenir dans
la sélection des coefficients soumis à comparaison.


\subsection{Hiérarchie de l’exposé
}

Le noyau formel est reproduit de manière commune avec les articles
précédents de la série. La répétition remplit une fonction d’autonomie :
chaque texte doit pouvoir être lu sans disposer des autres. Les extensions
spécifiques du présent travail apparaissent après le noyau et sont
identifiées par les objets qu’elles construisent.


Les coordonnées engendrées, la constante de structure fine et la constante
réduite de Planck interviennent lorsque les équations massiques les
requièrent. La paire angulaire précède les canaux et leurs moments ; les
registres précèdent la signature ; les fibres précèdent les opérateurs qui
agissent sur elles. L’électron central est conservé comme construction
structurelle propre, même lorsque sa grandeur est exprimée dans une échelle
chronologique partagée.


Dans cette hiérarchie, sélection, transport et mise à jour agissent
sur un même état enrichi. La composition de \(108\) pas
produit l’état terminal ; sa projection sur le registre dodécaphasique,
l’évaluation des canaux \(90/120\), l’extraction signée
\(U_{12}^{\rm sgn}\), la transformée orbitale et la reconstruction
entière établissent la chaîne

\[
 X_{\rm seed}\longmapsto X_{\rm term}\longmapsto
 C_{\rm term}\longmapsto U_{\rm term}\longmapsto K.
\]
Le registre, ses \(12\) composantes, ses quotients et sa mémoire proviennent
de cette évolution. Les coordinatisations des canaux, de la signature et de la période conservent dans un
même domaine le feuillet, l’orientation, la retenue et la frontière requis
par l’évaluation de masse. La composition~\eqref{vi:dep:cadena-terminal-completa}
réunit le producteur et sa reconstruction, avec les preuves de mise à jour et
d’inversion contiguës ; l’identité~\eqref{vi:dep:K-reversible} précise la
transformation réversible et son sous-réseau entier.


La définition de la particule prépare la classification. Sont ensuite
développés la conjugaison et l’holonomie, l’atlas, les observables
internes et l’évaluation de masse. La distribution par collections utilise
cette architecture, de sorte que chaque secteur est présenté avec son support,
son opération et sa lecture. Kepler--Sommerfeld se situe dans le développement
orbital de l’action, où peuvent être spécifiées les applications et les
conditions qui justifient la relation ; il n’est pas introduit comme une
ressemblance historique sans contenu opératoire.


\subsection{La comparaison comme résultat d’une chaîne déclarée
}

La comparaison requiert plus que deux colonnes de nombres. Il faut savoir ce
qui a été évalué, dans quelle unité et sous quelle définition de l’observable. Une
masse de pôle, une masse courante à une échelle déclarée et une limite
expérimentale ne sont pas interchangeables. Une largeur exprime une information
dynamique distincte de la position d’une masse. Le volume de l’inventaire
est utile lorsqu’il conserve ces distinctions, parce qu’il permet d’examiner
des collections complètes et pas seulement des exemples favorables.


Pour cette raison, la présentation réunit chaque sortie avec sa provenance et
avec les conditions de sa réalisation. Une identité interne,
l’évaluation d’une signature spécifiée et l’identification de cette signature
à une espèce sont des affirmations liées qui possèdent des démonstrations
propres. Les exposer consécutivement permet de montrer ce que détermine le
modèle à chaque niveau et quel est le contenu concret de la comparaison.


L’unité argumentative de l’article est, en définitive, la continuité entre
persistance, classification, opération et mesure. La masse est une sortie
de cette continuité ; la structure qui la produit reste disponible
pour comprendre ses relations avec les autres lectures de la particule.


% BEGIN SINTESIS_ADITIVA_20260922
\par
La section~\ref{delta22:vi:composicion} incorpore la longueur structurelle
avant l’horloge : l’incidence marquée produit \(N=5760\),
\(r_N=5759/23040\) et
\(L_\alpha=\alpha_{\rm HMT}^{16}r_NL_\star\)
dans~\eqref{delta22:vi:Lalpha} ; alors seulement sont fixés \(t_0\),
\(R_{108}=L_\alpha\), \(Q_L\), \(m_{\rm clk}\) et
\(b_e(Q_L)\). La loi de cotransformation
\eqref{delta22:vi:be-cotransformacion} montre pourquoi le facteur
électronique ne peut pas être figé lorsqu’on fait varier la section longitudinale ou
l’action. Les relations opératorielles \eqref{delta22:vi:R01} et
\eqref{delta22:vi:R02} éliminent la longueur d’action et déterminent le
coefficient gravitationnel unique~\eqref{delta22:vi:G-unica} ; la formule
circulaire~\eqref{delta22:vi:G-circular} s’obtient ensuite. La masse, son
caractère de route et ses tableaux restent causalement antérieurs à cette
identification radiale.

% END SINTESIS_ADITIVA_20260922
